% Motor: XeLaTeX. Sesión 26 - material docente UTP.
% Estructura docente de la Sesión 25.
\documentclass[aspectratio=169,10pt,c]{beamer}
% Reutiliza sin cambios la identidad visual UTP validada en la Sesion 08.
% Reutiliza sin cambios la identidad visual UTP validada en la Sesion 05.
% La segunda ruta permite compartir los assets comunes sin duplicar binarios.
% !TEX root = slides.tex
% Preambulo comun: estilo UTP, colores, fuentes, plantillas de beamer y
% comandos de apoyo didactico. Se carga desde slides.tex tras \documentclass.

% La opción "c" de beamer reparte el espacio libre 1:1.5 (arriba:abajo), por lo
% que el contenido queda por encima del centro. Igualamos ambos rellenos.
\makeatletter
\define@key{beamerframe}{c}[true]{%
  \beamer@frametopskip=0pt plus 1fill\relax%
  \beamer@framebottomskip=0pt plus 1fill\relax%
  \beamer@frametopskipautobreak=0pt plus .5\paperheight\relax%
  \beamer@framebottomskipautobreak=0pt plus .5\paperheight\relax%
}
\makeatother

\usepackage{graphicx}
\usepackage{tikz}
\usepackage{xcolor}
\usepackage{iftex}

\ifPDFTeX
  \usepackage[utf8]{inputenc}
  \usepackage[T1]{fontenc}
  \usepackage{lmodern}          % tipografia escalable para el respaldo con pdflatex
  \renewcommand{\familydefault}{\sfdefault}
\else
  \usepackage{fontspec}
  \IfFontExistsTF{Arial}
    {\setsansfont{Arial}[BoldFont={Arial Bold}]}
    {\setsansfont{Helvetica}[BoldFont={Helvetica Bold}]}
\fi

\usepackage[spanish,es-nodecimaldot]{babel}
\usepackage{microtype}
\usepackage{booktabs}
\usepackage{array}
\usetikzlibrary{arrows.meta,positioning,fit,calc,shapes.geometric}

\graphicspath{{assets/}}

% ---------- Identidad visual UTP ----------
\definecolor{utpRojo}{HTML}{C42B35}
\definecolor{utpNegro}{HTML}{111111}
\definecolor{utpGris}{HTML}{5F5F5F}
\definecolor{utpGrisMedio}{HTML}{9A9A9A}
\definecolor{utpGrisClaro}{HTML}{F1F1F1}

\usetheme{default}
\setbeamercolor{normal text}{fg=utpNegro,bg=white}
\setbeamercolor{structure}{fg=utpRojo}
\setbeamercolor{frametitle}{fg=utpNegro,bg=white}
\setbeamercolor{title}{fg=utpNegro}
\setbeamercolor{subtitle}{fg=utpNegro}
\setbeamercolor{author}{fg=utpNegro}
\setbeamercolor{date}{fg=utpNegro}
\setbeamercolor{block title}{fg=white,bg=utpNegro}
\setbeamercolor{block body}{fg=utpNegro,bg=utpGrisClaro}
\setbeamercolor{alertblock title}{fg=white,bg=utpRojo}
\setbeamercolor{alertblock body}{fg=utpNegro,bg=utpRojo!7}
\setbeamercolor{exampleblock title}{fg=white,bg=utpGris}
\setbeamercolor{exampleblock body}{fg=utpNegro,bg=utpGris!8}
\setbeamercolor{terminal}{fg=white,bg=utpNegro}

\setbeamerfont{title}{series=\bfseries,size=\fontsize{19.5}{22.5}\selectfont}
\setbeamerfont{subtitle}{size=\fontsize{11.5}{13.5}\selectfont}
\setbeamerfont{date}{series=\bfseries,size=\fontsize{11.5}{13.5}\selectfont}
\setbeamerfont{author}{size=\fontsize{10.5}{12.5}\selectfont}
\setbeamerfont{frametitle}{series=\bfseries,size=\large}
\setbeamerfont{block title}{series=\bfseries,size=\small}
\setbeamerfont{block body}{size=\normalsize}

\setbeamertemplate{navigation symbols}{}
\setbeamertemplate{headline}{}
\setbeamertemplate{blocks}[default]
\setbeamertemplate{itemize item}{\color{utpRojo}\small\raise.15ex\hbox{$\blacktriangleright$}}
\setbeamertemplate{itemize subitem}{\color{utpGris}\scriptsize$\bullet$}
\setbeamertemplate{enumerate item}{\color{utpRojo}\insertenumlabel.}
\setbeamertemplate{background canvas}{\color{white}\rule{\paperwidth}{\paperheight}}

% Título de diapositiva: negrita negra a la izquierda, logo UTP a la derecha
\setbeamertemplate{frametitle}{%
  \begin{tikzpicture}[remember picture,overlay]
    \node[anchor=north east,inner sep=0pt]
      at ([xshift=-7mm,yshift=-5.5mm]current page.north east)
      {\includegraphics[width=17mm]{utp-logo.png}};
  \end{tikzpicture}%
  \vspace{1.2mm}%
  \begin{minipage}[t]{.80\textwidth}\raggedright\insertframetitle\end{minipage}\par
  \vspace{1.1mm}{\color{utpRojo}\rule{\textwidth}{1.1pt}}
}

% Solo numeración discreta; la plantilla UTP no lleva pie de página
\setbeamertemplate{footline}{%
  \hfill{\color{utpGrisMedio}\scriptsize\insertframenumber}\hspace{6mm}\vspace{2.5mm}%
}

% Portada UTP
\setbeamertemplate{title page}{%
  \begin{tikzpicture}[remember picture,overlay]
    \node[anchor=center,inner sep=0pt] at (current page.center) {%
      \includegraphics[width=\paperwidth,height=\paperheight]{utp-cover-bg.jpeg}%
    };
    \node[anchor=north,align=center,text=utpNegro,inner sep=0pt,
      text width=145mm,font=\usebeamerfont{title}]
      at ([yshift=-25.2mm]current page.north) {\inserttitle};
    \node[anchor=north,align=center,text=utpNegro,inner sep=0pt,
      text width=142mm,font=\usebeamerfont{subtitle}]
      at ([yshift=-38.2mm]current page.north) {\insertsubtitle};
    \node[anchor=north,align=center,text=utpNegro,inner sep=0pt,
      font=\usebeamerfont{date}]
      at ([yshift=-50.4mm]current page.north) {\insertdate};
    \node[anchor=north,align=center,text=utpNegro,inner sep=0pt,
      font=\usebeamerfont{author}]
      at ([yshift=-58.4mm]current page.north) {\insertauthor};
  \end{tikzpicture}%
}

% ---------- Comandos de apoyo didáctico ----------
\newcommand{\concepto}[2]{%
  \begin{block}{#1}\large #2\end{block}}
\newcommand{\pregunta}[1]{%
  \begin{alertblock}{Decide y explica}#1\end{alertblock}}
\newcommand{\analogia}[1]{%
  \begin{exampleblock}{En palabras simples}\small #1\end{exampleblock}}
\newcommand{\enlarealidad}[1]{%
  \begin{block}{Cómo se ve en la realidad}\small #1\end{block}}
\newcommand{\ojo}[1]{%
  \begin{alertblock}{Ojo con esto}\small #1\end{alertblock}}
\newcommand{\apoyo}{{\normalfont\footnotesize\color{utpGris}~· material de apoyo}}
\newcommand{\marca}[1]{\tikz[baseline=(m.base)]{%
  \node[rounded corners=1.5pt,fill=utpGrisClaro,draw=utpGrisMedio,inner xsep=3pt,inner ysep=1.5pt,
        font=\footnotesize] (m) {#1};}}
\newenvironment{terminal}{%
  \begin{beamercolorbox}[wd=\linewidth,sep=5pt,rounded=true]{terminal}%
  \ttfamily\small\ignorespaces}{%
  \end{beamercolorbox}}

\tikzset{
  caja/.style={draw=utpNegro, very thick, rounded corners=1.5pt, fill=white,
    align=center, minimum height=0.72cm, inner sep=5pt, font=\small},
  cajaazul/.style={caja, draw=utpGris, fill=white},
  cajaverde/.style={caja, draw=utpNegro, fill=utpGrisClaro},
  cajanaranja/.style={caja, draw=utpRojo!75, fill=utpRojo!7},
  flecha/.style={-{Stealth[length=2.3mm]}, very thick, draw=utpNegro},
  flechavuelta/.style={{Stealth[length=2.3mm]}-{Stealth[length=2.3mm]}, very thick, draw=utpNegro}
}

\graphicspath{{../../Sesion05/slides/assets/}}



\graphicspath{{assets/}{../../Sesion05/slides/assets/}}
\usepackage{colortbl}
\usepackage{listings}
\lstset{basicstyle=\ttfamily\footnotesize,backgroundcolor=\color{utpGrisClaro},
  frame=none,aboveskip=4pt,belowskip=4pt,showstringspaces=false,
  columns=fullflexible,keepspaces=true,breaklines=true,
  keywordstyle=\color{utpRojo}\bfseries,
  morekeywords={FROM,WORKDIR,COPY,CMD},
  commentstyle=\color{utpGris}\itshape,morecomment=[l]{\#},
  xleftmargin=5pt,xrightmargin=5pt,framesep=5pt}

\title{Integración de Servicios Cloud GCP}
\subtitle{Balanceo de carga, CDN y contenedores}
\date{Semana 13 - Sesión 26}
\author{Rel Guzman}

\newcommand{\separador}[2]{%
\begin{frame}[plain]
\vfill
\begin{center}
{\fontsize{29}{34}\selectfont\bfseries #1}\par
\vspace{2mm}
{\small #2}\par
\vspace{4mm}
{\color{utpRojo}\rule{8.5cm}{1.2pt}}
\end{center}
\vfill
\end{frame}}

% Íconos simples para los diagramas.
\tikzset{
 caja/.style={draw=utpRojo,thick,rounded corners=4pt,align=center},
 vm/.style={draw=utpRojo,thick,rounded corners=3pt,fill=white,minimum width=1.6cm,minimum height=0.75cm,align=center,font=\scriptsize},
 vmmal/.style={vm,draw=utpGrisMedio,fill=utpGrisClaro,text=utpGris},
 zona/.style={draw=utpGrisMedio,dashed,rounded corners=4pt},
 flecha/.style={-latex,thick},
}
\newcommand{\cruz}[1]{\draw[utpRojo,line width=2pt] ($(#1)+(-0.25,-0.25)$) -- ($(#1)+(0.25,0.25)$) ($(#1)+(-0.25,0.25)$) -- ($(#1)+(0.25,-0.25)$);}

\begin{document}

{\setbeamertemplate{footline}{}\maketitle}

% 2
\begin{frame}{Revisión: lo que vimos en la Sesión 25}
\begin{columns}[T,onlytextwidth]
\column{.46\textwidth}
\begin{block}{Supervisar}
\small Mirar métricas y registros, y recibir una alerta cuando algo sale de lo normal.
\end{block}
\begin{block}{Alta disponibilidad}
\small Tener varias copias en zonas distintas, para que una falla no apague el servicio.
\end{block}
\column{.50\textwidth}
\begin{alertblock}{Responde}
\small
1. ¿Qué servicio de Google Cloud te envía un correo si la CPU de \texttt{web-1} pasa del 80\,\%?\par\medskip
2. ¿Qué es un punto único de falla?\par\medskip
3. ¿Por qué conviene tener VMs en zonas distintas?
\end{alertblock}
\end{columns}
\end{frame}

% 3
\begin{frame}{Revisión: la pregunta de hoy}
\begin{center}
\begin{tikzpicture}[font=\small]
  \foreach \i in {0,1,2} \node[caja,fill=utpRojo!12,minimum width=1.6cm] (c\i) at (0,1.1-\i*1.1) {cliente};
  \node[caja,dashed,minimum width=1.6cm,minimum height=1.2cm,font=\Large\bfseries] (q) at (3.6,0) {?};
  \foreach \i in {0,1,2} \draw[flecha] (c\i.east) -- (q.west);
  \foreach \n/\y in {1/1.1,2/0,3/-1.1} {
    \node[vm] (w\n) at (7.0,\y) {\texttt{web-\n}};
    \draw[flecha,dashed] (q.east) -- (w\n.west);
  }
\end{tikzpicture}
\end{center}
\vspace{3mm}
\begin{columns}[T,onlytextwidth]
\column{.32\textwidth}
\begin{block}{1. Balanceo de carga}
\small Tenemos 3 servidores. ¿Cómo se reparten los clientes?
\end{block}
\column{.32\textwidth}
\begin{block}{2. CDN}
\small ¿Cómo hacer que las fotos carguen rápido en Lima?
\end{block}
\column{.32\textwidth}
\begin{block}{3. Contenedores}
\small ¿Cómo poner el mismo programa en los 3 servidores sin instalar todo a mano?
\end{block}
\end{columns}
\end{frame}

% 4
\begin{frame}[c]
\begin{block}{Logro de la sesión}
Al finalizar, podrás explicar con tus palabras qué es un \textbf{balanceador de carga}, qué es una \textbf{CDN} y para qué sirven \textbf{Docker}, \textbf{Docker Compose} y \textbf{Kubernetes}. Además, verás en tu computadora cómo Nginx reparte clientes entre tres servidores.
\end{block}
\end{frame}

% 5
\begin{frame}{Saberes previos}
\begin{itemize}\setlength\itemsep{14pt}
  \item Si 100 personas abren a la vez una página que corre en un solo servidor, ¿qué crees que pasa?
  \item Desde Lima abres la página de una tienda de EE.\,UU. ¿Crees que carga igual de rápido que una tienda con servidores en Lima?
  \item Un compañero te pasa su proyecto y en tu computadora no funciona. ¿Qué podría faltar?
\end{itemize}
\end{frame}

% 6
\begin{frame}[c]{Utilidad de esta sesión}
\begin{center}
\begin{minipage}{.84\textwidth}
{\large Cuando una aplicación tiene muchos clientes, se ponen \textbf{varias copias} de ella. Hoy verás cómo se reparten los clientes entre esas copias, cómo se entregan las fotos desde cerca y cómo se encienden copias iguales con un solo comando.}
\end{minipage}
\end{center}
\end{frame}

% 7
\separador{Parte 1: balanceo de carga}{Repartir a los clientes entre varios servidores}

% 8
\begin{frame}{El problema: un solo servidor}
\begin{block}{El caso}
\small Seis clientes piden un reporte de ventas al mismo tiempo. Cada reporte tarda \textbf{5 segundos} y el servidor hace \textbf{uno a la vez}.
\end{block}
\vspace{6mm}
\begin{center}
\begin{tikzpicture}[font=\small]
  \node[vm,minimum width=2.0cm,minimum height=1.1cm,font=\small] (s) at (0,0) {\textbf{servidor}};
  \foreach \i/\t in {1/5,2/10,3/15,4/20,5/25,6/30} {
    \node[caja,fill=utpRojo!\the\numexpr 6+\i*4\relax,minimum width=1.35cm,font=\scriptsize] (c\i) at (0.6+\i*1.65,0) {cliente \i};
    \node[font=\scriptsize,anchor=north] at (c\i.south) {espera \textbf{\t\,s}};
  }
  \draw[flecha] (c1.west) -- (s.east);
\end{tikzpicture}
\end{center}
\vspace{6mm}
\begin{alertblock}{Se forma una cola}
\small El cliente 6 espera \textbf{30 segundos}. ¿Y si hubiera 3 servidores?
\end{alertblock}
\end{frame}

% 9
\begin{frame}{La solución: el balanceador de carga}
\begin{center}
\begin{tikzpicture}[font=\small]
  \foreach \i in {0,...,3} \node[caja,fill=utpRojo!12,minimum width=1.5cm,font=\footnotesize] (c\i) at (0,1.65-\i*1.1) {cliente};
  \node[caja,fill=utpRojo!30,text width=2.8cm,minimum height=1.3cm] (lb) at (4.3,0) {\textbf{Balanceador}\\\footnotesize una sola dirección};
  \foreach \i in {0,...,3} \draw[flecha] (c\i.east) -- (lb.west);
  \foreach \n/\y in {1/1.2,2/0,3/-1.2} {
    \node[vm] (w\n) at (8.6,\y) {\texttt{web-\n}};
    \draw[flecha] (lb.east) -- (w\n.west);
  }
\end{tikzpicture}
\end{center}
\vspace{2mm}
\begin{block}{Balanceador de carga}
\small Recibe a todos los clientes en \textbf{una sola dirección} y los \textbf{reparte} entre varios servidores iguales.
\end{block}
\begin{columns}[T,onlytextwidth]
\column{.48\textwidth}
\begin{exampleblock}{El cliente}
\small Solo conoce una dirección. No sabe cuántos servidores hay detrás.
\end{exampleblock}
\column{.48\textwidth}
\begin{exampleblock}{Dónde lo verás}
\small En Google Cloud: \textbf{Cloud Load Balancing}. En la práctica de hoy: \textbf{Nginx}.
\end{exampleblock}
\end{columns}
\end{frame}

% 10
\begin{frame}{Round robin: repartir por turnos}
\begin{columns}[c,onlytextwidth]
\column{.40\textwidth}
\begin{block}{Round robin}
\small Un cliente para cada servidor, en orden, y se vuelve a empezar.
\end{block}
\renewcommand{\arraystretch}{1.15}
\centering\small
\begin{tabular}{@{}cl@{}}
\toprule
\textbf{Cliente} & \textbf{Va a}\\
\midrule
1 & \texttt{web-1}\\
2 & \texttt{web-2}\\
3 & \texttt{web-3}\\
\rowcolor{utpRojo!20} 4 & \texttt{web-1} {\color{utpGris}\scriptsize(otra vuelta)}\\
5 & \texttt{web-2}\\
6 & \texttt{web-3}\\
\bottomrule
\end{tabular}
\column{.56\textwidth}
\centering{\color{utpRojo}\bfseries Los 6 reportes de antes}\par\vspace{4mm}
\begin{tikzpicture}[font=\small,x=0.17cm]
  \node[anchor=east,font=\footnotesize] at (0,0) {1 servidor};
  \fill[utpGrisMedio] (0.5,-0.3) rectangle (30.5,0.3);
  \node[anchor=west,font=\footnotesize\bfseries] at (31,0) {30 s};
  \node[anchor=east,font=\footnotesize] at (0,-1.1) {3 servidores};
  \fill[utpRojo] (0.5,-1.4) rectangle (10.5,-0.8);
  \node[anchor=west,font=\footnotesize\bfseries] at (11,-1.1) {10 s};
\end{tikzpicture}
\par\vspace{5mm}
\begin{alertblock}{Por qué}
\small Cada servidor hace 2 reportes, y los tres trabajan \textbf{al mismo tiempo}.
\end{alertblock}
\end{columns}
\end{frame}

% 11
\begin{frame}{¿Y si un servidor se cae?}
\begin{center}
\begin{tikzpicture}[font=\small]
  \node[caja,fill=utpRojo!12,minimum width=1.6cm] (c) at (0,0) {clientes};
  \node[caja,fill=utpRojo!30,minimum width=2.4cm,minimum height=1.1cm] (lb) at (3.6,0) {\textbf{Balanceador}};
  \draw[flecha] (c) -- (lb);
  \node[vm] (w1) at (8.0,1.2) {\texttt{web-1}};
  \node[vmmal] (w2) at (8.0,0) {\texttt{web-2}};
  \node[vm] (w3) at (8.0,-1.2) {\texttt{web-3}};
  \draw[flecha] (lb.east) -- (w1.west);
  \draw[thick,dashed,utpGrisMedio] (lb.east) -- (w2.west);
  \draw[flecha] (lb.east) -- (w3.west);
  \cruz{w2}
\end{tikzpicture}
\end{center}
\vspace{3mm}
\begin{block}{Verificación de estado}
\small Cada pocos segundos, el balanceador le pregunta a cada servidor si responde. Si uno no responde, \textbf{deja de enviarle clientes} y reparte entre los que quedan.
\end{block}
\begin{alertblock}{Lo que ve el cliente}
\small Nada raro: la página carga igual. Cuando \texttt{web-2} vuelve, el balanceador lo usa otra vez.
\end{alertblock}
\end{frame}

% 12
\separador{Parte 2: CDN}{Copias de los archivos cerca del cliente}

% 13
\begin{frame}{¿Qué es una CDN?}
\begin{block}{CDN (red de entrega de contenido)}
\small Servidores en muchas ciudades del mundo que guardan una \textbf{copia} de tus archivos (fotos, logo) y se los entregan al cliente \textbf{desde cerca}.
\end{block}
\vspace{2mm}
\begin{center}
\begin{tikzpicture}[font=\small]
  \node[caja,fill=utpRojo!12,minimum width=1.8cm,minimum height=0.9cm] (u) at (0,0) {cliente\\\footnotesize en Lima};
  \node[caja,fill=utpRojo!30,text width=2.6cm,minimum height=1.1cm] (cdn) at (4.6,0) {\textbf{Copia}\\\footnotesize cerca de Lima};
  \node[caja,fill=white,text width=2.6cm,minimum height=1.1cm] (vm) at (10.4,0) {\textbf{Tus servidores}\\\footnotesize en EE.\,UU.};
  \draw[flecha] (u.east) -- node[above,font=\scriptsize]{\texttt{logo.png}} (cdn.west);
  \draw[flecha,dashed] (cdn.east) -- node[above,font=\scriptsize]{solo la primera vez} (vm.west);
\end{tikzpicture}
\end{center}
\vspace{2mm}
\begin{columns}[T,onlytextwidth]
\column{.48\textwidth}
\begin{block}{La primera vez}
\small La CDN no tiene el logo: lo pide a tus servidores, se lo entrega al cliente y \textbf{guarda una copia}.
\end{block}
\column{.48\textwidth}
\begin{block}{Las siguientes veces}
\small La CDN entrega su copia. Es más rápido y tus servidores \textbf{ni se enteran}.
\end{block}
\end{columns}
\end{frame}

% 14
\begin{frame}{¿Qué va en la CDN?}
\renewcommand{\arraystretch}{1.35}
\centering\small
\begin{tabular}{@{}>{\raggedright\arraybackslash}p{5.6cm}>{\raggedright\arraybackslash}p{5.6cm}@{}}
\toprule
\textbf{Sí va en la CDN} & \textbf{No va en la CDN}\\
\midrule
Logo & El carrito de Ana\\
Fotos de productos & Los pedidos de Luis\\
Archivos CSS y JavaScript & El estado de un pago\\
\bottomrule
\end{tabular}
\par\vspace{5mm}
\begin{columns}[T,onlytextwidth]
\column{.48\textwidth}
\begin{alertblock}{La regla}
\small Va lo que es \textbf{igual para todos}. No va lo que \textbf{cambia según el cliente}: Luis podría recibir la copia del carrito de Ana.
\end{alertblock}
\column{.48\textwidth}
\begin{block}{En Google Cloud}
\small Se llama \textbf{Cloud CDN} y se activa con una casilla en el balanceador de carga.
\end{block}
\end{columns}
\end{frame}

% 15
\separador{Parte 3: contenedores}{Docker, Docker Compose y Kubernetes}

% 16
\begin{frame}{El problema: preparar cada servidor}
\begin{columns}[T,onlytextwidth]
\column{.48\textwidth}
\begin{block}{Para la práctica de hoy necesitas}
\small
1. Instalar Node.js.\par\smallskip
2. Instalar Nginx, que se instala distinto en Windows, macOS y Linux.\par\smallskip
3. Abrir cinco terminales.
\end{block}
\column{.48\textwidth}
\begin{alertblock}{Ahora imagina}
\small Hacerlo en 3 servidores. Si en uno instalas otra versión de Node.js, ese servidor podría fallar.
\end{alertblock}
\end{columns}
\par\vspace{5mm}
\begin{block}{Contenedor (lo viste en la Sesión 12)}
\small Un paquete con tu programa \textbf{y todo lo que necesita} para funcionar, para que corra igual en cualquier máquina.
\end{block}
\end{frame}

% 17
\begin{frame}{Docker: imagen y contenedor}
\begin{center}
\begin{tikzpicture}[font=\small]
  \node[caja,fill=white,text width=2.3cm,minimum height=1.4cm] (df) at (0,0) {\textbf{Dockerfile}\\\footnotesize la receta};
  \node[caja,fill=utpRojo!30,text width=2.3cm,minimum height=1.4cm] (img) at (4.6,0) {\textbf{Imagen}\\\footnotesize el paquete};
  \draw[flecha] (df) -- node[above,font=\scriptsize]{crear} (img);
  \foreach \n/\y in {1/1.1,2/0,3/-1.1} {
    \node[vm,minimum width=2.2cm] (k\n) at (9.6,\y) {contenedor \n};
    \draw[flecha] (img.east) -- (k\n.west);
  }
  \node[font=\scriptsize] at (6.9,1.0) {encender};
\end{tikzpicture}
\end{center}
\vspace{3mm}
\begin{columns}[T,onlytextwidth]
\column{.48\textwidth}
\begin{block}{Imagen}
\small El paquete guardado: el programa y todo lo que necesita. Está apagado.
\end{block}
\column{.48\textwidth}
\begin{block}{Contenedor}
\small Una imagen \textbf{encendida}. De una imagen puedes encender muchos contenedores iguales.
\end{block}
\end{columns}
\par\vspace{3mm}
\raggedright\small \textbf{Docker} es el programa que crea las imágenes y enciende los contenedores.
\end{frame}

% 18
\begin{frame}[fragile]{El Dockerfile de la tienda}
\begin{columns}[c,onlytextwidth]
\column{.44\textwidth}
\begin{lstlisting}
FROM node:22-alpine
WORKDIR /app
COPY package.json server.js ./
CMD ["node", "server.js", "3001"]
\end{lstlisting}
\column{.52\textwidth}
\renewcommand{\arraystretch}{1.3}
\small
\begin{tabular}{@{}>{\ttfamily}l>{\raggedright\arraybackslash}p{4.4cm}@{}}
\toprule
\textsf{\textbf{Línea}} & \textbf{En palabras}\\
\midrule
FROM & Empieza con Node.js ya instalado\\
WORKDIR & Trabaja en la carpeta \texttt{/app}\\
COPY & Copia los archivos de la tienda\\
CMD & Al encender, ejecuta la tienda\\
\bottomrule
\end{tabular}
\end{columns}
\par\vspace{5mm}
\begin{lstlisting}
docker build -t tienda .          # crea la imagen "tienda"
docker run -p 3001:3001 tienda    # enciende un contenedor
\end{lstlisting}
\par\vspace{2mm}
\begin{exampleblock}{Lo importante}
\small No instalaste Node.js en tu computadora: viene dentro de la imagen.
\end{exampleblock}
\end{frame}

% 19
\begin{frame}{Docker Compose: varios contenedores a la vez}
\begin{center}
\begin{tikzpicture}[font=\small]
  \node[caja,fill=white,text width=2.4cm,minimum height=1.4cm] (f) at (0,0) {\textbf{compose.yaml}\\\footnotesize un archivo};
  \draw[utpRojo,thick,rounded corners=4pt] (3.4,-1.3) rectangle (10.6,1.4);
  \node[font=\scriptsize\bfseries,color=utpRojo,anchor=south west] at (3.4,1.4) {docker compose up};
  \node[caja,fill=utpRojo!30,minimum width=1.6cm,minimum height=0.9cm] (n) at (4.7,0) {\textbf{nginx}};
  \foreach \n/\y in {1/0.9,2/0,3/-0.9} {
    \node[vm,minimum width=1.8cm,minimum height=0.6cm] (t\n) at (8.9,\y) {\texttt{tienda\n}};
    \draw[flecha] (n.east) -- (t\n.west);
  }
  \draw[flecha] (f) -- (3.4,0);
\end{tikzpicture}
\end{center}
\begin{block}{Docker Compose}
\small Describe en \textbf{un archivo} todos los contenedores de tu aplicación y los enciende juntos con \textbf{un comando}.
\end{block}
\renewcommand{\arraystretch}{1.1}
\centering\small
\begin{tabular}{@{}>{\raggedright\arraybackslash}p{5.6cm}>{\raggedright\arraybackslash}p{5.6cm}@{}}
\toprule
\textbf{A mano (Práctica 1)} & \textbf{Con Docker Compose (Práctica 2)}\\
\midrule
Instalar Node.js y Nginx & Solo tener Docker\\
5 terminales, 5 comandos & 1 comando: \texttt{docker compose up}\\
\bottomrule
\end{tabular}
\end{frame}

% 20
\begin{frame}{Kubernetes: contenedores en muchas máquinas}
\begin{columns}[c,onlytextwidth]
\column{.50\textwidth}
\centering
\begin{tikzpicture}[font=\small]
  \draw[utpRojo,thick,rounded corners=4pt] (-0.15,-1.15) rectangle (5.95,1.55);
  \foreach \n/\x in {1/0.9,2/2.9,3/4.9} {
    \draw[zona] (\x-0.85,-0.85) rectangle (\x+0.85,1.3);
    \node[font=\scriptsize,color=utpGris,anchor=north] at (\x,1.3) {máquina \n};
  }
  \node[vm,minimum width=1.3cm,minimum height=0.5cm] (p1) at (0.9,0.35) {tienda};
  \node[vmmal,minimum width=1.3cm,minimum height=0.5cm] (p2) at (2.9,0.35) {tienda};
  \cruz{p2}
  \node[vm,minimum width=1.3cm,minimum height=0.5cm] (p3) at (4.9,0.35) {tienda};
  \node[vm,minimum width=1.3cm,minimum height=0.5cm,dashed] (p4) at (4.9,-0.45) {tienda nueva};
  \draw[flecha,dashed] (p2.south) to[bend right=25] (p4.west);
\end{tikzpicture}
\column{.46\textwidth}
\begin{block}{Kubernetes}
\small Tú le dices: «quiero \textbf{3 copias} de la tienda». Kubernetes las reparte en varias máquinas y se encarga de que \textbf{siempre haya 3}.
\end{block}
\begin{exampleblock}{Si se cae una copia}
\small Kubernetes lo nota y enciende otra, sin que tú hagas nada.
\end{exampleblock}
\end{columns}
\par\vspace{5mm}
\begin{columns}[T,onlytextwidth]
\column{.48\textwidth}
\begin{alertblock}{¿Por qué no basta Compose?}
\small Compose usa \textbf{una sola máquina}. Si esa máquina se apaga, se apaga todo.
\end{alertblock}
\column{.48\textwidth}
\begin{block}{En Google Cloud}
\small Se llama \textbf{Google Kubernetes Engine} (GKE).
\end{block}
\end{columns}
\end{frame}

% 21
\begin{frame}{Las tres herramientas}
\renewcommand{\arraystretch}{1.5}
\centering\small
\begin{tabular}{@{}p{2.4cm}>{\raggedright\arraybackslash}p{3.3cm}>{\raggedright\arraybackslash}p{3.3cm}>{\raggedright\arraybackslash}p{3.3cm}@{}}
\toprule
 & \textbf{Docker} & \textbf{Docker Compose} & \textbf{Kubernetes}\\
\midrule
\textbf{Qué hace} & Crea la imagen y enciende un contenedor & Enciende varios contenedores juntos & Mantiene muchas copias encendidas\\
\textbf{Dónde} & Una máquina & Una máquina & Varias máquinas\\
\textbf{Ejemplo} & La imagen de la tienda & Práctica 2: 3 tiendas + Nginx & La tienda en producción\\
\bottomrule
\end{tabular}
\par\vspace{6mm}
\begin{alertblock}{En una frase}
\small Docker \textbf{empaqueta}, Docker Compose \textbf{junta} y Kubernetes \textbf{mantiene}.
\end{alertblock}
\end{frame}

% 22
\begin{frame}[fragile]{Práctica 1: balanceo con Nginx en tu computadora}
\begin{block}{Caso de estudio}
\small Pedidos Express prepara su tienda online para el Cyber Day: tres copias del servidor (Node.js) detrás de \textbf{Nginx}, que reparte por turnos: 3001, 3002, 3003, 3001\dots{} (\textit{round robin}).
\end{block}
\begin{columns}[c,onlytextwidth]
\column{.38\textwidth}
\begin{tikzpicture}[font=\small]
  \node[caja,fill=utpRojo!12,minimum width=1.3cm,font=\footnotesize] (c) at (0,0) {clientes};
  \node[caja,fill=utpRojo!30,minimum width=1.3cm,align=center] (n) at (2.0,0) {\textbf{Nginx}\\\footnotesize :8080};
  \draw[flecha] (c) -- (n);
  \foreach \p/\y in {3001/0.9,3002/0,3003/-0.9} {
    \node[vm,minimum width=1.1cm,minimum height=0.6cm] (s\p) at (4.0,\y) {\texttt{\p}};
    \draw[flecha] (n.east) -- (s\p.west);
  }
\end{tikzpicture}
\column{.60\textwidth}
\renewcommand{\arraystretch}{1.15}
\scriptsize
\begin{tabular}{@{}lll@{}}
\toprule
\textbf{Ruta} & \textbf{Qué hace} & \textbf{El servidor\dots}\\
\midrule
\texttt{GET /} & Página principal & responde al instante\\
\texttt{GET /productos} & Consulta la base de datos & \textbf{espera} 2 s\\
\texttt{POST /comprar} & Cobra con el banco & \textbf{espera} 5 s\\
\texttt{GET /reporte} & Calcula un reporte & \textbf{trabaja} 5 s con la CPU\\
\bottomrule
\end{tabular}
\end{columns}
\par\vspace{3mm}
\begin{exampleblock}{Cómo empezar (carpeta \texttt{Sesion26/practica\_balanceo/})}
\footnotesize
\begin{tabular}{@{}l>{\ttfamily}l@{}}
1. Tres terminales, un servidor en cada una: & node server.js 3001 \quad\textsf{(y 3002, 3003)}\\
2. Otra terminal, para encender Nginx: & nginx -p \string"\$PWD\string" -c nginx.conf\\
3. Otra terminal, para simular clientes: & node prueba.js / 6 --en-fila\\
\end{tabular}
\end{exampleblock}
\end{frame}

% 23
\begin{frame}{Práctica 1: seis pruebas}
\renewcommand{\arraystretch}{1.12}
\centering\small
\begin{tabular}{@{}c>{\ttfamily\scriptsize\raggedright\arraybackslash}p{5.5cm}>{\raggedright\arraybackslash}p{6.6cm}@{}}
\toprule
\textbf{\#} & \textsf{\small\textbf{Comando}} & \textbf{Qué observar}\\
\midrule
1 & node prueba.js / 6 --en-fila & ¿En qué orden responden los servidores?\\
2 & node prueba.js /productos 6 --directo 3001 & Un solo servidor: ¿cuánto tarda en total?\\
3 & node prueba.js /reporte 6 --directo 3001 & Un solo servidor: ¿cuánto espera el sexto cliente?\\
4 & node prueba.js /reporte 6 & Con Nginx: ¿cuánto tarda en total?\\
5 & node prueba.js /reporte 1\newline \textsf{\footnotesize y, antes de 5 s, en otra terminal:}\newline node prueba.js / 6 & ¿Qué visitas tardan? ¿A qué servidor le tocaron?\\
6 & \textsf{\footnotesize Apaga el 3002 (}Ctrl+C\textsf{\footnotesize ) y luego:}\newline node prueba.js / 6 --en-fila & ¿Algún cliente ve un error? ¿Qué servidores responden?\\
\bottomrule
\end{tabular}
\par\vspace{3mm}
\begin{alertblock}{La idea}
\small Cuando el servidor \textbf{espera}, atiende a muchos a la vez. Cuando \textbf{trabaja} con la CPU, los clientes hacen cola: ahí ayuda repartir entre varios servidores.
\end{alertblock}
\end{frame}


% 24
\begin{frame}[fragile]{Práctica 2: la misma tienda con Docker Compose}
\begin{columns}[T,onlytextwidth]
\column{.50\textwidth}
\begin{block}{Antes de empezar}
\small Apaga los servidores y Nginx de la Práctica 1. Necesitas \textbf{Docker Desktop}. Si no lo tienes, sigue la demostración del docente.
\end{block}
\begin{lstlisting}
docker compose up --build -d
docker compose ps
node prueba.js / 6 --en-fila
docker compose down
\end{lstlisting}
\column{.46\textwidth}
\begin{exampleblock}{Qué observar}
\small
1. ¿Cuántos contenedores muestra \texttt{docker compose ps}?\par\medskip
2. ¿La prueba da lo mismo que en la Práctica 1?\par\medskip
3. ¿Cuántos comandos usaste para encender todo? ¿Y en la Práctica 1?
\end{exampleblock}
\end{columns}
\par\vspace{3mm}
{\footnotesize\color{utpGris} Los archivos \texttt{Dockerfile} y \texttt{compose.yaml} están en la misma carpeta. La guía completa está en el \texttt{README.md}.}
\end{frame}

% 25
\begin{frame}{Cierre: preguntas para comprobar}
\begin{block}{Responde con tus palabras}
\small
1. ¿Qué hace un balanceador de carga?\\\medskip
2. Con 3 servidores y round robin, ¿a cuál va el cliente 4?\\\medskip
3. Da un ejemplo de archivo que sí va en una CDN y uno que no.\\\medskip
4. ¿Qué diferencia hay entre una imagen y un contenedor?\\\medskip
5. ¿Qué hace Kubernetes que Docker Compose no hace?
\end{block}
\begin{alertblock}{Idea final}
\small Muchos clientes se atienden con varias copias iguales: el balanceador las reparte, la CDN entrega las fotos desde cerca y los contenedores hacen que todas las copias sean idénticas.
\end{alertblock}
\end{frame}

% 26
\begin{frame}{Referencias}
\scriptsize
\begin{itemize}\setlength\itemsep{4pt}
  \item Prácticas: \texttt{github.com/rgap/rgap-Integracion-de-Servicios-Cloud-Gcp-Agosto2026}, carpeta \texttt{Sesion26/practica\_balanceo/}.
  \item Nginx: \texttt{nginx.org}. Node.js: \texttt{nodejs.org}.
  \item Docker y Docker Compose: \texttt{docs.docker.com}. Kubernetes: \texttt{kubernetes.io}.
  \item Google Cloud: Cloud Load Balancing, Cloud CDN y Google Kubernetes Engine (GKE): \texttt{cloud.google.com}.
  \item Material previo del curso: Sesión 12 (contenedores) y Sesión 25 (supervisión y alta disponibilidad).
  \item Las respuestas de las preguntas están en \texttt{../respuestas.md}.
\end{itemize}
\end{frame}

\end{document}
